\chapter{Fehlerdiagnose und optionale PETSc-Umgebung}\label{sec:app-diagnosis}
\index{Fehlerdiagnose}\index{PETSc}
\section{Systematisch vom Eingang zur Bilanz prüfen}
\begin{longtable}{p{.26\linewidth}p{.64\linewidth}}
\toprule Beobachtung & Ursache prüfen und nächste Handlung\\\midrule\endhead
Import schlägt fehl & Unterstützte Python-Version, verwendeten Interpreter,
\path{pip check} und Installation desselben Wheels prüfen.
\path{python -I} mit absolutem Beispielpfad ausführen; kein
\path{PYTHONPATH}-Workaround als Artefaktnachweis.\\
Modellpfad oder Schema ungültig & Dokumentpfad der \path{ModelSchemaError}
lesen, TOML-Schlüssel gegen das zutreffende Schema vergleichen,
relative Dateien innerhalb des Modellverzeichnisses halten.
Keine stillschweigende Versionsmigration.\\
Material nicht zulässig & Einheiten, Tensorform, Symmetrie/Positivität,
Voigt-Reihenfolge und lokale Achsen kontrollieren. Negative Leitfähigkeit
ist kein gültiger Testparameter für einen erfolgreichen Fall.\\
Geometrie abgelehnt & Degenerierte Zellen, Außenflächenbesitz,
Tet10-Kantenlage und Gmsh-Auswahlmengen prüfen; mm/m-Skala explizit setzen.\\
Nullraum/Singularität & Komponenten einzeln kontrollieren: Wärmeanker,
sechs mechanische Starrkörpermoden, elektrische Gauge. Schwebende
Elektrode fixiert keine Referenz. Keine künstliche Stabilisierung, die
unbemerkt das Problem verändert.\\
Iterativer Solver konvergiert nicht & Originalresiduen je Feld,
Randwerte, Konditionierung und Skalierung zuerst prüfen. Erst danach
bewusst Iterationsbudget/Präkonditionierung ändern; Methode im Bericht
festhalten. Kein automatischer SciPy-Fallback behauptet.\\
Bilanz oder Energie falsch & Einheiten/Normalenvorzeichen,
Volumen-/Flächenlast und zeitliche Speicherterme vergleichen.
Relative Residuen ersetzen keine Bilanzprüfung.\\
Ausgabeziel existiert & Bei normalen Exporten einen neuen Ordner oder
Präfix wählen. Nur Study verwendet dasselbe Ziel zur kontrollierten
Wiederaufnahme. Keine fremden Ergebnisse löschen.\\
Abgebrochener Lauf & Study-Katalog erneut prüfen und denselben Auftrag
fortsetzen. Zeitserienstatus \path{incomplete} bezeichnet lesbare Teilframes,
keinen automatischen Neustart eines Zeitsolvers.\\
Bericht/Surrogat inkompatibel & Schema, Einheiten, CSV-/VTU-/JSON-Hashes
und Kandidatenversion abgleichen. Unvollständige Ausgabe nicht konsumieren;
mit ursprünglichen Eingaben in neues Ziel neu erzeugen.\\
Surrogat extrapoliert & Parameterbox und Spaltenreihenfolge kontrollieren.
Außerhalb neue FEM-Fälle erzeugen oder ausdrücklich als Extrapolation
kennzeichnen; keine allgemeine Unsicherheitszusage.\\\bottomrule
\caption{Fehlerhilfe ohne unbemerkte Veränderung des physikalischen Modells.}
\end{longtable}
Fehlerbericht enthält minimalen Fall, Inputdateien/Hashes, Paket-/Python-
und Dependency-Versionen, tatsächlichen Importpfad, erwarteten SI-Wert,
Istwert und vorab definierte Toleranz. Private Daten und Zugangsdaten
gehören nicht in öffentliche Issues. Bild und Solver-Status allein genügen
zur Reproduktion nicht.

\section{Native PETSc-Auswahl nur über die direkte API}
\path{PetscBackend()} wählt PREONLY/LU, \path{PetscFieldSplitBackend()}
FGMRES mit Rechtspräkonditionierung und FieldSplit/Schur für zwei Felder.
Beide werden explizit mit \path{backend=...} an direkte Solver übergeben.
\path{backend=None} bleibt SciPy; TOML-Modelle und Studien konfigurieren
PETSc derzeit nicht. Die acht bestehenden Backend-Hooks ersetzen weder
Mathematik noch Randbedingungen. Der Eigenwertlöser bleibt SciPy; PETSc kann
bei Moden die elektrische lineare Elimination übernehmen.

Geprüfte native Referenz ist der im Repository definierte Ubuntu-24.04-
Container mit Python 3.12, PETSc/petsc4py 3.19.6 und getrennten
\path{float64}/\path{complex128}-Runtimes. Ein Prozess verwendet
\path{COMM_SELF}; MPI-Verteilung und Performancegewinne werden nicht
zugesagt. Ein generischer pip-Installationsversuch ohne Scalar-/ABI-Vertrag
ist keine verifizierte PETSc-Umgebung. Auf einem vollständigen Checkout:
\begin{lstlisting}
docker build -f ci/petsc.Dockerfile -t fem-toolbox-petsc:test .
docker run --rm -v "$PWD:/work:ro" fem-toolbox-petsc:test python ci/run_petsc_tests.py --scalar real -q tests/petsc -p no:cacheprovider
docker run --rm -v "$PWD:/work:ro" fem-toolbox-petsc:test python ci/run_petsc_tests.py --scalar complex -q tests/petsc -p no:cacheprovider
\end{lstlisting}
Diese Infrastruktur ist im Repository vorhanden, nicht Bestandteil des
Handbuch-Beispielarchivs. Die CI führt beide nativen Varianten separat aus.
Vor dem ersten Import wird ein Scalar-Runtime gewählt; ihn nicht im laufenden
Prozess wechseln. Ein reeller Build verarbeitet komplexe Systeme über die
verifizierte Realblockabbildung, ein komplexer direkt. Kein Imaginärteil
wird verworfen. Fehlende oder inkompatible Installation wird erst beim
Erzeugen des Backends diagnostiziert; der Core importiert ohne PETSc.

FGMRES erhält explizite \path{rtol}, \path{atol},
\path{max_iterations} und \path{restart}. Mehr als zwei nichtleere Felder
sind nicht unterstützt; ein Feld nutzt native LU-Präkonditionierung.
Methoden-/Konvergenzbericht und unskalierte Feldresiduen prüfen.
Ein direkter LU-Bericht besitzt keine erfundene Iterationszahl.
Nichtkonvergenz ist ein Fehler, keine stillschweigende Backend-Umschaltung.

\section{Nachweise und Pflege des Benutzerbands}
Die maschinenlesbare Workflowmatrix \path{recipes.json} verbindet jedes
Programm mit Befehl, Eingaben, Ausgabe, Sollwert, Grenzen und Fehlerhilfe.
Die vollständigen Originaldateien liegen im Quellenarchiv.
\path{tests/unit/test_user_handbook_sources.py} prüft Inventar, Byteidentität,
Hashes und Einbindung aller Programme im PDF. Bei einer Beispieländerung
werden Kopie, Rezepttext und Nachweis im selben PR aktualisiert.
Die nachvollziehbare Kandidatenprüfung steht in
\path{Doc/latex/USER_VERIFICATION.md}; sie nennt ausgeführte Befehle,
Versionen, isolierte Artefaktumgebung, Toleranzen und Prüflücken.
Der Linux-Wheel-/sdist-CI-Lauf ist ein zusätzlicher unabhängiger Runtime,
keine erfundene zweite Person. Redaktionelle Eigenprüfung ersetzt keine
unabhängige Freigabe. API-/Entwicklungsband \#209, Nachweisband \#210 und
Gesamtabnahme \#211 bleiben eigene Arbeitspakete.
